\documentclass[11pt]{article}
\usepackage[a4paper,margin=1in]{geometry}
\usepackage{amsmath,amsfonts,amssymb}
\usepackage{graphicx}
\usepackage{caption}
\usepackage[hidelinks]{hyperref}
\usepackage{microtype}

\setlength{\textfloatsep}{8pt plus 2pt minus 2pt}
\setlength{\floatsep}{6pt plus 2pt minus 2pt}
\setlength{\intextsep}{6pt plus 2pt minus 2pt}
\setlength{\abovecaptionskip}{4pt}
\setlength{\belowcaptionskip}{0pt}

\begin{document}

\noindent Throughout this section, $\kappa$ denotes the membrane permeability, $T$ the
target-hit time, $\tau_{\mathrm{out}}$ the first exit time from the corridor band, and
$\tau_{\mathrm{mem}}$ the first genuine membrane crossing. The three windows
\texttt{peak1}, \texttt{valley}, and \texttt{peak2} are compared using two complementary
descriptions: \emph{normalized} observables, reported as ratios or percentages of $T$, and
\emph{absolute} budgets, reported in expected step counts. More specifically, for an event
time $\tau$, the relative timing ratio is
\[
  \frac{E[\tau \mid \tau < T,\ T \in W]}{E[T \mid T \in W]},
\]
i.e., the mean event time expressed as a fraction of the mean hit time in window $W$;
the outside-time share is the fraction of $T$ spent outside the corridor band among
trajectories whose hit times fall in the corresponding window; the outside budget is the
expected number of pre-hit steps spent outside the corridor band; and the post-crossing
budget is the expected number of steps remaining after the first membrane crossing and
before target hit. In Figure~\ref{fig:fig1-geometry-curve-budget-en}, ``outside'' is
the union of the left-side pocket and the merged outer/right-side bulk.

\begin{figure}[!htbp]
  \centering
  \caption{Top: spatial partition for the symmetric one-target baseline. Bottom: hit-time
    curves $f(t)$ for $\kappa=0$ and $\kappa=0.0040$, with embedded stacked bars for the
    \texttt{peak1}, \texttt{valley}, and \texttt{peak2} windows marked on the curves. The
    three segments give the normalized time shares in the left side, corridor, and merged
    outer/right side. The $t_1$, $t_v$, and $t_2$ markers indicate the true window
    centres; the bars are shifted horizontally only for legibility.}
  \label{fig:fig1-geometry-curve-budget-en}
  \includegraphics[width=0.95\linewidth]{figures/fig1_geometry_curve_budget}
\end{figure}

Figure~\ref{fig:fig1-geometry-curve-budget-en} first locates the late-time structure in
space. The curve separates the direct \texttt{peak1} arrivals from a valley and a later
\texttt{peak2} window, but this visual feature is not labelled \texttt{double\_peak}
unless the quantitative classifier supports that label; weaker cases should remain
shoulders or local bumps. The budget bars show why the later window is delayed. The
outside share rises from about $48\%$ to $56\%$ at $\kappa=0$ and from about $47\%$ to
$56\%$ at $\kappa=0.0040$, with the left-side pocket reaching about $30\%$ in the later
window. Thus the second-peak window is not just a plotting artefact: it is enriched in
paths that spend extra pre-hit time outside the corridor, especially in the left-side
pocket, before returning to the target. This additional residence time supplies the
extra time scale behind the late-time feature.

\clearpage

\begin{figure}[!htbp]
  \centering
  \caption{Relative timing ratio for $\tau_{\mathrm{out}}$, defined as
    $E[\tau_{\mathrm{out}} \mid \tau_{\mathrm{out}} < T,\ T \in W]/E[T \mid T \in W]$,
    that is, the mean exit time expressed as a fraction of the mean hit time in window
    $W$. Black labels give the relative timing ratios, colour-matched percentages give
    the outside-time shares, and the top-right annotation gives the corresponding outside
    budgets, namely the expected numbers of pre-hit steps spent outside the corridor
    band, for the valley and \texttt{peak2} windows. The three series correspond to
    $\kappa=0$, $\kappa=0.0040$, and the boundary value $\kappa=0.0152$.}
  \label{fig:fig2-tau-out-budget-en}
  \includegraphics[width=0.85\linewidth]{figures/fig2_tau_out_budget}
\end{figure}

Figure~\ref{fig:fig2-tau-out-budget-en} shows that the timing of the first corridor exit
is only part of the story. \texttt{peak1} behaves as the direct-corridor benchmark: its
exit ratio is about $0.76$--$0.79$ and its outside-time share is only $9\%$--$10\%$.
The valley and \texttt{peak2} windows, however, both have small exit ratios, so the
normalized exit time alone does not explain why one window lies later. The separation
appears in the absolute outside budget: it increases from $718$ to $1311$ steps at
$\kappa=0$ and from $705$ to $1288$ steps at $\kappa=0.0040$. The mechanism is therefore
not simply ``early exit'', but long residence after exit. Together with Figure~1, this
points to outside excursions, especially through the left-side pocket, as the source of
the delayed mass in the full first-passage distribution rather than a statement about
the mean first-passage time alone.

\clearpage

\begin{figure}[!htbp]
  \centering
  \caption{Relative timing ratio for $\tau_{\mathrm{mem}}$, i.e., the mean first-crossing
    time expressed as a fraction of the mean hit time in window $W$. Black labels give the
    relative timing ratios, colour-matched percentages give the membrane-crossing
    probabilities $P(\tau_{\mathrm{mem}} < T \mid T \in W)$, and the top-right annotation
    gives the post-crossing budgets, namely the expected numbers of steps remaining after
    the first crossing and before target hit. The label `n/a' denotes that the ratio is
    undefined when no crossing occurs before target hit. The three series correspond to
    $\kappa=0$, $\kappa=0.0040$, and the boundary value $\kappa=0.0152$.}
  \label{fig:fig3-tau-mem-budget-en}
  \includegraphics[width=0.85\linewidth]{figures/fig3_tau_mem_budget}
\end{figure}

Figure~\ref{fig:fig3-tau-mem-budget-en} separates the membrane-mediated part of the same
delay mechanism. When $\kappa=0$, crossing and post-crossing budgets vanish, as they
should. Once $\kappa>0$, the valley and \texttt{peak2} windows are still close in
crossing probability, about $18\%$ versus $20\%$ at $\kappa=0.0040$, but they differ
more clearly in the time left after the first crossing: about $213$ versus $336$ steps,
increasing to about $604$ and $785$ steps at $\kappa=0.0152$. Permeability therefore
does not act mainly by making crossing more frequent. Its visible effect is to enlarge
the post-crossing budget, so the trajectories that enter this channel still have a
longer remaining path to absorption. This supports a delayed-residence explanation for
the shoulder or second-peak contribution, while any \texttt{double\_peak} label remains
a classifier-level claim.

\end{document}
